\section{代表性研究与案例分析}

\subsection{光伏预测：深度学习已经成为主流路线}
Mellit 等人在光伏输出功率预测综述中指出，光伏预测已经从传统统计方法逐步转向机器学习与深度学习，并且随着多源气象数据与更高频采样数据的可获得性提升，深度学习的重要性持续增强\cite{mellit2020}。该文强调，光伏预测并不是单纯的时间序列回归问题，而是与天气、季节、地理位置和时间尺度紧密相关的综合建模任务。这个结论说明，模型选择不能脱离具体应用背景。

Khouili 等人在 2025 年的系统综述中进一步指出，在纳入分析的 26 篇深度学习光伏预测研究中，LSTM 是使用频率最高的架构，占比 32.69\%，CNN 紧随其后，占比 28.85\%\cite{khouili2025}。同时，多源数据融合、可解释性和鲁棒性仍然是未来研究重点。由此可以看出，深度学习已经成为光伏预测的主流技术路线，但工程化问题依然没有完全解决。

\subsection{光伏案例一：LSTM、BiLSTM 与 CNN-LSTM 的横向比较}
Mellit、Massi Pavan 和 Lughi 基于意大利 Trieste 大学微电网的真实光伏数据，对 LSTM、BiLSTM、GRU、BiGRU、CNN1D、CNN-LSTM 和 CNN-GRU 等网络进行了系统比较\cite{mellit2021}。该研究同时覆盖 1 分钟、5 分钟、30 分钟和 60 分钟等多种时间尺度，因此非常适合说明不同模型在实际光伏场景中的性能差异。研究表明，在一步预测且时间尺度较短时，结构相对简单的 LSTM 或 GRU 已经可以取得较高精度；而当任务转为多步预测、天气更复杂时，模型性能差异会被明显放大。

这个案例说明深度学习在短时光伏预测中确实具有优势，但预测步长越长，问题就越困难。它的意义不仅在于证明“某个模型效果更好”，更在于提供了统一数据源、多模型、多时间尺度比较的研究框架，这比只展示单次最优结果更有说服力。

\subsection{光伏案例二：领域知识融合比单纯堆模型更重要}
Luo、Zhang 和 Zhu 提出的 PC-LSTM 为光伏预测提供了很有启发性的思路\cite{luo2021}。该研究尝试把领域知识和物理约束引入深度学习框架，以提升小时级和日前光伏功率预测的合理性与稳定性。作者指出，在训练数据稀疏或场景变化明显时，纯数据驱动方法容易得到“误差看起来不高、但在物理上并不合理”的结果，而加入领域知识的模型往往表现更稳健。

这一思路对智慧能源具有普遍意义。真实系统中的很多问题都不是“数据越多、网络越深”就一定更好。气象异常、传感器故障和站点迁移都会导致纯黑箱模型性能下降。因此，把物理约束、设备边界或经验规律融入模型，是深度学习在能源场景中真正走向可用的重要方向。

\subsection{光伏案例三：面向泛化能力的元学习方法}
Sarmas 等人在 2023 年提出了一种基于元学习的短时光伏预测方法，用多个 LSTM 变体作为基础模型，再通过元学习器动态组合预测结果\cite{sarmas2023}。该方法在葡萄牙里斯本的三套屋顶光伏系统上进行了验证，并且不依赖数值天气预报输入。研究结果显示，该方法相比单一最佳基础模型可带来最高约 5\% 的精度提升，相比简单平均集成也有进一步改善。

这个案例的重要性在于，它说明深度学习研究的重点正在从“谁的单模型更强”逐渐转向“如何提升跨场景泛化与稳定性”。对智慧能源这类强场景依赖问题来说，泛化能力往往比单站点上的最低误差更重要。

\subsection{风电案例：深度学习正从时序建模走向时空建模}
尽管本文主线聚焦光伏预测，但风电预测同样能反映深度学习在智慧能源中的发展趋势。Liu 和 Zhang 在 2024 年的综述中系统梳理了 140 余篇深度学习风电预测研究，指出风电预测正在从传统历史功率或单点气象输入，逐渐过渡到结合 SCADA、数值天气预报和图结构信息的更复杂建模框架\cite{liu2024}。这说明新能源预测正在从单序列建模走向时空协同建模。

更具体地说，Meka、Alaeddini 和 Bhaganagar 基于一个 130 MW、包含 86 台风机的真实风场数据，利用 TCN 模型对 0 至 50 分钟的短时风电功率进行多步预测，并与 LSTM、CNN-LSTM 和多元线性回归模型进行比较\cite{meka2021}。结果表明，TCN 在不同风速区间下都表现出更稳定的优势。这个案例可以作为光伏研究的横向补充，说明深度学习在“高频、波动、强时序依赖”的新能源任务上具有普遍适用性。

\subsection{微电网调度案例：预测价值最终要落到决策上}
若只讨论预测误差，容易忽略智慧能源的最终目标是更优运行。Alabdullah 和 Abido 在微电网能量管理研究中采用 Deep Q-Network 构建强化学习调度框架，并在真实微电网案例上验证了方法有效性\cite{alabdullah2022}。研究表明，该方法在随机负荷、分布式电源和电价波动条件下能够获得接近最优调度的结果，其运行成本与最优调度方案相比相差在 2\% 以内。

这一研究提醒我们，预测模型的真正价值不在于单独优化某个误差指标，而在于为储能充放电、购电决策、负荷转移和系统安全留出更好的决策空间。因此，在写作课程报告时，最好把“预测”与“调度”联系起来，而不是把它写成一个孤立任务。

\begin{center}
\captionof{table}{代表性论文案例对比\\Table 2 Comparison of representative studies}
\vspace{3pt}
{\fontsize{7.8pt}{11pt}\selectfont
\setlength{\tabcolsep}{2pt}
\begin{tabular}{@{}p{0.17\columnwidth}p{0.19\columnwidth}p{0.18\columnwidth}p{0.20\columnwidth}p{0.18\columnwidth}@{}}
\toprule
文献 & 场景与数据 & 模型/方法 & 关键发现 & 可支撑的观点 \\
\midrule
Mellit 等，2021\cite{mellit2021} & Trieste 大学微电网光伏数据；1、5、30、60 分钟尺度 & LSTM、BiLSTM、GRU、CNN-LSTM 等 & 一步预测精度高，多步预测仍具挑战 & 深度学习适合短时光伏预测，但预测步长越长越难 \\
Luo 等，2021\cite{luo2021} & 小时级、日前光伏功率预测 & PC-LSTM & 融合领域知识后更稳健，稀疏数据下更有优势 & 机理与数据融合优于纯黑箱建模 \\
Khouili 等，2025\cite{khouili2025} & 26 篇深度学习光伏预测研究综述 & 系统综述 & LSTM 与 CNN 最常见，鲁棒性与解释性仍是挑战 & 深度学习已成主流，但工程化问题仍突出 \\
Liu 和 Zhang，2024\cite{liu2024} & 140 余篇风电预测研究综述 & 系统综述 & 风电预测从时序模型走向时空与图结构建模 & 深度学习在新能源预测中具有跨场景适用性 \\
Meka 等，2021\cite{meka2021} & 130 MW 风场、86 台风机、12 个月数据 & TCN vs. LSTM/CNN-LSTM/MLR & TCN 在多步短时预测中更稳定 & 高频多步预测需要更强时序建模能力 \\
Alabdullah 和 Abido，2022\cite{alabdullah2022} & 真实微电网能量管理案例 & DQN 强化学习 & 运行成本与最优调度相差在 2\% 以内 & 预测最终要服务于调度优化 \\
Sarmas 等，2023\cite{sarmas2023} & 葡萄牙三套屋顶光伏系统 & 元学习 + 多 LSTM & 相比单一最佳基础模型可提升约 5\% & 泛化能力和动态集成是重要方向 \\
\bottomrule
\end{tabular}}
\label{tab:case-comparison}
\end{center}
